\SourceLeaf{324}%
or, if one substitutes the value from (3) for \(u\) and uses (8), (9), and (11) to set
\[
s=\frac{r\cos\theta^2}{a\cos\varphi}=\frac{r\cos\theta\operatorname{tang}P}{a\operatorname{tang}Q}
\]
one obtains
\[
t=\frac{r\cos\theta\operatorname{tang}P\mathord{.}\sin U}{a}(1+\tfrac1{24}uu+\tfrac1{12}ss-\frac{ee\cos P^2}{6\cos\varphi^2}\mathord{.}qq)
\]
For \(\cos\theta\mathord{.}\operatorname{tang}P=\BareRoot{(1-ee\sin P^2)}\mathord{.}\operatorname{tang}P\), equation (4) gives the value expanded to the order needed here
\[
\BareRoot{(1-ee\sin B^2)}\mathord{.}\operatorname{tang}B(1+\frac{3ee-6e^4\sin P^2+3e^4\sin P^4}{8\cos\varphi^2(1-ee\sin P^2)}\mathord{.}qq)
\]
and consequently
\[
\begin{aligned}
t={}&\frac{rk\operatorname{tang}B\sin U}{a}(1+\frac{5ee+(4ee-14e^4)\sin P^2+5e^4\sin P^4}{24\cos\varphi^2(1-ee\sin P^2)}\mathord{.}qq\\
&\qquad+\tfrac1{12}ss+\tfrac1{24}uu)
\end{aligned}
\]
If one now makes the same substitutions here as in the preceding article, the final result is (III)
\[
\begin{aligned}
t={}&\frac{rk\operatorname{tang}B\mathord{.}\sin T}{a}(1+\frac{5ee+(4ee-14e^4)\sin B^2+5e^4\sin B^4}{24k^4}\mathord{.}bb\\
&\qquad+\frac{k^4}{12aa\cos\varphi^2}\mathord{.}rr+\tfrac1{24}tt)
\end{aligned}
\]

\indent The three formulas I, II, III give, in essence, the solution to our problem. Their derivation shows directly that they are accurate through the third order inclusive. That their accuracy in fact extends one order further, or that the error in each of the formulas is of fifth order, could easily be shown by a few supplementary intermediate expansions. It could also be shown from the fact that the expressions, by their nature, cannot contain quantities of even order. I will pass over this point, however, since the \textit{second} derivation of the formulas, to be carried out in the following articles (26---32), yields the same result on its own.

\Article{26}
This investigation is to be treated as self-contained and independent of everything that precedes it. For convenience and to avoid uncertainty, all the notation used in it will therefore be explained as it is introduced. Most letters already used in the first derivation will retain the meanings they had there. Two of them (\(u\) and \(s\)), however, since they represent only auxiliary quantities%
